% Section 1.10: Complete Project Example
\section{Complete Project Example: From Zero to Flashable Binary}
\label{sec:complete-project-example}

Now that we've covered the individual components, let's walk through creating a complete, flashable embedded Rust project from scratch. This section provides a comprehensive example that you can follow step-by-step.

\subsection{Project Overview}
\label{subsec:project-overview}

We'll create a minimal but complete embedded Rust project with all necessary files:

\begin{itemize}
    \item Project structure and configuration
    \item Cargo configuration for cross-compilation
    \item Build script for linker integration
    \item Memory layout specification
    \item Minimal working application
\end{itemize}

\subsection{Project Structure}
\label{subsec:project-structure}

A complete embedded Rust project requires several files beyond just \filepath{main.rs}. Here's the full structure:

\begin{minted}{text}
efr32mg24-minimal/
├── .cargo/
│   └── config.toml      # Build configuration
├── src/
│   ├── main.rs          # Entry point
│   └── lib.rs           # Library root (if needed)
├── build.rs             # Build script for linker
├── memory.x             # Memory layout for linker
├── Cargo.toml           # Package metadata
├── rust-toolchain.toml  # Rust version/toolchain
└── .gitignore
\end{minted}
\subsection{Cargo Configuration (\.cargo/config.toml)}
\label{subsec:cargo-config}
This file is critical for embedded builds. It tells Cargo:

\begin{itemize}
    \item Which target to compile for
    \item What runner to use (probe-rs for flashing/debugging)
    \item Environment variables for debugging
\end{itemize}

\begin{minted}{toml}
[build]
target = "thumbv8m.main-none-eabihf"  # ARM Cortex-M33 target

[target.thumbv8m.main-none-eabihf]
runner = "probe-rs run --chip EFR32MG24B220F1536IM48"

[env]
DEFMT_LOG = "debug"
\end{minted}
